\section{Estrategia de seguridad, permisos y aislamiento}

Los perfiles funcionales son Operaciones y Logística, Comercial y Compras, y
Administración y Calidad. Cada actor pertenece a un solo perfil. El permiso relaciona
perfil y clase documental; si la combinación no existe, el acceso se deniega.
El nivel de sensibilidad describe el dato, pero no reemplaza esta matriz.

Los roles técnicos separan tareas. \path{rag_ingesta} incorpora borradores y
vectores; \path{rag_documental} publica o revoca mediante funciones atómicas;
\path{rag_runtime} consulta y registra interacciones; y \path{rag_auditor}
reconstruye evidencia en modo lectura. Todos son
\texttt{NOLOGIN}, sin privilegios de superusuario ni \texttt{BYPASSRLS}. La
autenticación queda para una aplicación futura y fuera del TP.

Cada transacción fija el código del actor en \texttt{rag.actor}. Las funciones
de contexto comprueban que exista, esté activo y tenga un perfil activo. Un
valor ausente, vacío, inexistente o inactivo produce un contexto nulo y las
políticas deniegan la lectura. La base deriva el perfil desde el actor para que
los valores no se contradigan. No autentica quién fija \texttt{rag.actor}: una
aplicación futura deberá asociar la identidad autenticada con el actor de la
transacción.

RLS se aplica y se fuerza sobre \texttt{documento},
\texttt{version\_documental}, \texttt{fragmento} y \texttt{embedding}. También
limita condiciones comerciales e incidencias según el perfil. La vista de
recuperación usa los privilegios del invocador, no los del propietario. Las 15
combinaciones perfil--clase, los cuatro contextos inválidos y el acceso directo
a tablas se prueban desde un rol no propietario.

La auditoría es append-only para los roles operativos: no pueden insertar
eventos arbitrarios, modificarlos ni eliminarlos. Los triggers generan eventos
para consultas, respuestas y evidencias; las funciones documentales registran
publicación, sustitución y revocación dentro de la misma transacción. El
propietario conserva \texttt{TRUNCATE} para recrear el conjunto sintético, una
operación administrativa que no se concede a los roles de aplicación.
